\section{Ideals}

\begin{theorem}[Maximal ideal existence]
\label{mcrt-auto-theorem-1-1}\dcref{1.1}
\source{matsumura-commutative-ring-theory:page-0017}\uses{}
Every proper ideal $I$ of a ring $A$ is contained in a maximal ideal.
\end{theorem}
\begin{proof}
The set of proper ideals containing $I$, ordered by inclusion, is inductive:
the union of a chain is again a proper ideal.  Zorn's lemma therefore gives a
maximal element.
\end{proof}

\begin{definition}[Prime ideal]
An ideal $\mathfrak p\subsetneq A$ is prime when $A/\mathfrak p$ is an integral
domain, equivalently $xy\in\mathfrak p$ implies $x\in\mathfrak p$ or
$y\in\mathfrak p$.
\uses{}
\end{definition}

\begin{definition}[Multiplicative set]
A subset $S\subseteq A$ is multiplicative when $1\in S$ and $st\in S$ for all
$s,t\in S$.
\uses{}
\end{definition}

\begin{theorem}[Prime ideal disjoint from a multiplicative set]
\label{mcrt-auto-theorem-1-2}\dcref{1.2}
\source{matsumura-commutative-ring-theory:page-0017}\uses{}
If $I$ is an ideal disjoint from a multiplicative set $S$, then there exists a
prime ideal $\mathfrak p\supseteq I$ with $\mathfrak p\cap S=\varnothing$.
\end{theorem}
\begin{proof}
Choose by Zorn a maximal ideal $\mathfrak p$ among ideals containing $I$ and
disjoint from $S$.  If $x,y\notin\mathfrak p$, maximality gives
$s\in(\mathfrak p+xA)\cap S$ and $t\in(\mathfrak p+yA)\cap S$.  Their product
lies in $\mathfrak p+xyA$ and in $S$, so $xy\notin\mathfrak p$.
\end{proof}

\begin{definition}[Radical and Jacobson radical]
For an ideal $I$ put
\[
\sqrt I=\{a\in A:a^n\in I\text{ for some }n>0\},\qquad
\operatorname{nil}(A)=\sqrt{(0)}.
\]
The Jacobson radical is the intersection of all maximal ideals,
$\operatorname{rad}(A)=\bigcap_{\mathfrak m\in\operatorname{Max}A}\mathfrak m$.
A ring with one maximal ideal is local, and a ring with finitely many maximal
ideals is semilocal.
\uses{}
\end{definition}

\begin{theorem}[Pairwise comaximal ideals]
\label{mcrt-auto-theorem-1-3}\dcref{1.3}
\source{matsumura-commutative-ring-theory:page-0018}\uses{}
If $I_1,\ldots,I_r$ are pairwise comaximal ideals, then
\[
I_1\cap\cdots\cap I_r=I_1\cdots I_r.
\]
\end{theorem}
\begin{proof}
For two ideals, $I+J=A$ gives
$I\cap J=(I+J)IJ=IJ$.  Induct on $r$, observing that $I_1\cdots I_{r-1}$
is comaximal with $I_r$.
\end{proof}

\begin{theorem}[Chinese remainder theorem]
\label{mcrt-auto-theorem-1-4}\dcref{1.4}
\source{matsumura-commutative-ring-theory:page-0019}\uses{mcrt-auto-theorem-1-3}
If $I_1,\ldots,I_r$ are pairwise comaximal, the canonical map
\[
A/(I_1\cdots I_r)\longrightarrow\prod_{i=1}^r A/I_i
\]
is an isomorphism.
\end{theorem}
\begin{proof}
The kernel is $\bigcap_i I_i=\prod_i I_i$ by the preceding theorem.  For two
ideals choose $e\in I_1$ and $f\in I_2$ with $e+f=1$; then
$af+be$ maps to $(a,b)$.  Induction proves surjectivity for $r$ ideals.
\end{proof}

\section{Modules}
\begin{theorem}[Theorem 2.1]\label{mcrt-auto-theorem-2-1}\dcref{2.1}\source{matsumura-commutative-ring-theory:page-0022}\uses{}
\begin{mathematicalrecord}
Suppose that M is an A-module generated by n elements,
 and that cpEHom,(M, M); let I be an ideal of A such that q(M) c ZM.
 Then there is a relation of the form
    (**) (p~+alcp?-l+~~~+u,-,~+u,=O,
 with Ui~Z? for 1 d i 6 n (where both sides are considered as endomorph-
 isms of M).
\end{mathematicalrecord}
\end{theorem}
\begin{proof}
\begin{mathematicalrecord}
Let M = Ao, + ... + AU,; by the assumption cp(M) c IM there exist
 Uij~l such that cp(oi) = cJ= i uijwj. This can be rewritten

          jil   ((~6, - 6,) uj = 0 (for 1 d i < n),

where aij is the Kronecker symbol. The coefficients of this system of linear
equations can be viewed as a square matrix ((~6~~- aij) of elements of A?[q],
the commutative subring of the endomorphism ring E of M generated by
the image A? of A together with cp; let b, denote its (i,j)th cofactor, and
8         Commutative   rings and modules

d its determinant. By multiplying the above equation through by b, and
summing over i, we get do, = 0 for 1 d k d n. Hence d.M = 0, so that d = 0
as an element of E. Expanding the determinant d gives a relation of the
form (**). W
\end{mathematicalrecord}
\end{proof}


\begin{theorem}[Theorem 2.2]\label{mcrt-auto-theorem-2-2}\dcref{2.2}\source{matsumura-commutative-ring-theory:page-0023}\uses{}
\begin{mathematicalrecord}
NAK). Let M be a finite A-module and I an ideal of A. If
M = IM    then there exists aEA such that aM = 0 and a = 1 mod I. If in
addition I c rad (A) then M = 0.
\end{mathematicalrecord}
\end{theorem}
\begin{proof}
\begin{mathematicalrecord}
Setting cp= 1, in the previous theorem gives the relation a =
l+a,+    .. . + a, = 0 as endomorphisms   of M, that is aM = 0, and
a = 1 modI. If I c rad(A) then a is a unit of A, so that on multiplying
both sides of aM = 0 by a-l we get M L 0. n
\end{mathematicalrecord}
\end{proof}


\begin{theorem}[Theorem 2.3]\label{mcrt-auto-theorem-2-3}\dcref{2.3}\source{matsumura-commutative-ring-theory:page-0023}\uses{}
\begin{mathematicalrecord}
Let (A, m, k) be a local ring and M a finite A-module; set
R = M/mM. Now M is a finite-dimensional      vector space over k, and we
?2       Modules

write n for its dimension. Then:
   (i) If we take a basis {tii,..., U,) for M over k, and choose an
inverse image UiE M of each y then {ul ,. . . ,u} is a minimal basis
of M;
   (ii) conversely every minimal basis of M is obtained in this way, and
so has n elements.
   (iii) If {q,. . .,u} and {ui ,. . . ,u} are both minimal bases of M, and
ui = 1 UijUj with Uij~A then det (aij) is a unit of A, SO that (aij) is an
invertible matrix.
\end{mathematicalrecord}
\end{theorem}
\begin{proof}
\begin{mathematicalrecord}
(i) M = xAui + mM, and M is finitely generated (hence also
M/xAui),       so that by the above corollary M = xAui. If {ul,. .,u} is
not minimal, so that, for example, {u2,. . . , un} already generates M
then (&,..., tin} generates fi, which is a contradiction.                 Hence
{ul, _. . , u,> is a minimal basis.
   (ii) If {ui, . . . , u} is a minimal basis of M and we set Ui for the image
of ui in M, then u,,. . . , U, generate A, and are linearly independent
over k; indeed, otherwise some proper subset of {tii , . . . ,U} would
be a basis of M, and then by (i) a proper subset of {ui,. . ,u} would
generate M, which is a contradiction.
                                                                  - -
   (iii) Write Zij for the image in k of aij, so that Ui = Caijuj holds in
M. Since (aij) is the matrix transforming one basis of the vector space
I$ into another, its determinant is non-zero. Since det (aij) modm =
det (aij) # 0 it follows that det (aij) is a unit of A. By Cramer?s formula
the inverse matrix of (aij) exists as a matrix with entries in A. n
   We give another interesting application of NAK, the proof of which is
due to Vasconcelos [2].
\end{mathematicalrecord}
\end{proof}


\begin{theorem}[Theorem 2.4]\label{mcrt-auto-theorem-2-4}\dcref{2.4}\source{matsumura-commutative-ring-theory:page-0024}\uses{}
\begin{mathematicalrecord}
Let A be a ring and M a finite A-module. If f:M -M             is
an A-linear map and f is surjective then f is also injective, and is thus
an automorphism     of M.
\end{mathematicalrecord}
\end{theorem}
\begin{proof}
\begin{mathematicalrecord}
Since f commutes with scalar multiplication       by elements of A, we
can view M as an A[X]-module        by setting X.m = f(m) for meM. Then by
assumption XM = M, so that by NAK there exists YEA[X] such that
(1 +XY)M    =O. Now        for uEKer(f)        we have 0 = (1 + XY)(u) =
u + Yf(u) = u, so that f is injective.     n
\end{mathematicalrecord}
\end{proof}


\begin{theorem}[Theorem 2.5]\label{mcrt-auto-theorem-2-5}\dcref{2.5}\source{matsumura-commutative-ring-theory:page-0024}\uses{}
\begin{mathematicalrecord}
Let (A,m) be a local ring; then a projective module over A
is free (for the definition of projective module, see Appendix B, p. 277).
\end{mathematicalrecord}
\end{theorem}
\begin{proof}
\begin{mathematicalrecord}
This is easy when M is finite: choose a minimal basis ol,. . . , w, of
M and define a surjective map cp:F -+ M from the free module F =
Ae, O...@Ae, to M by cp(xaiei) = c a,~,,.. ?f1 we set K = Ker(cp) then, from
10       Commutative      rings and modules

the minimal     basis property,
         Capi     = 0 s- aiEm      for all i.
Thus K c mF. Because M is projective, there exists $: M -F     such that
F = $(M)@    K, and it follows that K = mK. On the other hand, K is a
quotient of F, therefore finite over A, so that K = 0 by NAK and F N M.
  The result was proved by Kaplansky [2] without the assumption that
M is finite. He proves first of all the following lemma, which holds for
any ring (possibly non-commutative).

Lemma 1. Let R be any ring, and F an R-module which is a direct sum of
countably generated submodules; if M is an arbitrary direct summand ofF
then M is also a direct sum of countably generated submodules.
Proof of Lemma 1. Suppose that F = M @ N, and that F = oloA E,, where
each E, is countably generated. By translinite induction, we construct a
well-ordered family {F} of submodules of F with the following properties:
      (i) ifa<Bthen        F,cF,,
     (ii) F = ua F,,
   (iii) if c1is a limiting ordinal then F, = UBcaFB,
   (iv) F,, ,/F, is countably generated,
     (v) F,= M,@N,,          where M,= MnF,,   N,= NnF,,
   (vi) each F, is a direct sum of E, taken over a suitable subset of A.
We now construct such a family {FE). Firstly, set F, = (0). For an
ordinal cc,assume that F, has been defined for all ordinals p < CI. If a is
a limiting ordinal, set F, = Us<= F,. If tl is of the form a=,&+ 1, let Q1
be any one of the E, not contained in F, (if F, = F then the construction
stops at Fp). Take a set xrl, x12,. . . of generators of Qi, and decompose
xi1 into its M- and N-components; now let Q, be the direct sum of the
finitely many E, which are necessary to write each of these two components
in the decomposition F = GE,, and let xzl, xz2,. . . be generators of Qz.
Next decompose xi2 into its M- and N-components, let Q3 be the direct
sum of the finitely many E, needed to write these components, and let
x3l,     x32,*..  be generators of Q3. Then carry out the same procedure with
xzl, getting xbl, xa2,. . . , then do the same for x13. Carrying out the same
procedure for each of the xij in the order x~l,x~2,xz1,x~3,x22,x31,...
we get ?countably many elements xij. We let F, be the submodule of F
generated by F, and the xii, and this satisfies$all our requirements. This
gives the family {Fb} .
      Now M = u M,, with each M, a direct summand of F, and M,, 1 =) M,,
so that M, is also a direct summand of M,+l. Moreover,
         F,+ JFn = W,+ ,lMJOW,+            JNa),
?2       Modules                                                                11

and hence M,, ,/M, is countably generated. Thus we can write
        M a+l = M,@M&+,,         with Mi,, countably generated.
When a is a limit ordinal, since M, = up <aM,, we set Mh = 0. Then finally
we can write
        M = @ Mi with Mb (at most) countably generated. n
   Of course aolfree module satisfies the assumption of Lemma 1, so that,
in particular, we see that any projective module is a direct sum of countably
generated projective modules. Thus in the proof of Theorem 2.5 we can
assume that M is countably generated.
Lemma 2. Let M be a projective module over a local ring A, and XEM.
Then there exists a direct summand of M containing x which is a free
module.
Proof of Lemma 2. We write M as a direct summand of a free module
F = M @ N. Choose a basis B = {ui}ipl of F such that the given element
x has the minimum         possible number of non-zero coordinates when
expressed in this basis. Then if x = ulal + ... + ~,a,, with 0 # Ui~A, we have
          ai4 1 Aaj for i = 1, 2,. . . , n;
             J#i
indeed, if, say, a, = 1; - i biai then x = C;- i (ui + u,b,)q, which contradicts
the choice of B. Now set ui = yi + zi with y,gM and z,EN; then
          X=&Ui=piyi.
If we write yi = I;= 1cijuj + ti, with ti linear combinations of elements of
B other than ui, . . . , u,, we get relations a, = CJ= I ajcji, and, hence, in view
of what we have seen above, we must have
          1 - ci+m       and cijEm for i #j.
It follows that the matrix (cij) has an inverse (this can be seen from the
fact that the determinant is s 1 mod m, or by elimination). Thus replacing
Ul,..., unbyy,,...,     y, in B, we still have a basis of F. Hence, F, = CYiA is
a direct summand of F, and hence also of M, and satisfies all the
requirements of Lemma 2. n
   To prove the theorem, let M be a countably generated projective module,
M=co,A+o,~+....              By Lemma 2, there exists a free module F, such
that o,EF,, and M = F, GM,, where M, is a projective module. Let w;
be the M ,-component of o2 in the decomposition M = F I @M 1, and take
a free module F, such that o;EF~ and M, = F, GM,, where M, is a
Projective module. Let oj be the M,-component                 of wj in M = F, @
F, 0 M,; proceeding in the same way, we get
         M=F,@F,@...,
so that M is a free module.       n
12       Commutative    rings and modules

   We say that an A-module M # 0 is a simple module if it has no
submodules other than 0 and M itself For any 0 # UEM, we then have
M = Ao. Now Ao N A/ann(o), but in order for this to be simple, arm(o)
must be a maximal ideal of A. Hence, any simple A-module is isomorphic
to A/m with m a maximal ideal, and conversely an A-module of this form
is simple. If M is an A-module, a chain

of submodules of M is called a composition series of M if every Mi/Mi+ 1 is
simple; r is called the length of the composition series. If a composition series
of M exists, its length is an invariant of M independent of the choice of
composition series. More precisely, if M has a composition series of length
r, and if MI N, I>... =3N, is a strictly descending chain of sub-
modules, then we have s < r. This invariance corresponds to part of the
basic JordanHolder        theorem in group theory, but it can easily be proved
on its own by induction, and the reader might like to do this as an exercise.
The length of a composition series of M is called the length of M, and written
2(M); if M does not have a composition series we set l(M) = co. A necessary
and sufficient condition for the existence of a composition series of M is that
the submodules of M should satisfy both the ascending and descending
chain conditions (for which see 93). In general, if N c M is a submodule,
we have
         l(M) = l(N) + l(M/N).
IfO+M,     -M2     -..?   - M, -+ 0 is an exact sequence of A-modules and
each Mi has finite length then
         i$l (- 1)?NMJ= 0.
   If m is a maximal ideal of A and is finitely generated over A then
1(A/m?) < co. In fact,
          1(A/m?) = 1(A/m) + l(m/m2) + ... + l(mv-l/mv);
now each m?/m i+l is a finite-dimensional        vector space over the field
k = A/m, and since its A-submodules are the same thing as its vector
subspaces, ,(mi/mifl)     is equal to the dimension of m?/m?+ l as k-vector
space. (This shows that A/m? is an Artinian ring, see $3.)
   Considering /(A/m?) for all v, we get a function of v which is intimately
related to the ring structure of A, and which also plays a role in the resolution
of singularities in algebraic and complex analytic geometry; this is studied in
Chapter 5.
   We say that an A-module M is of finite presentation if there exists an
exact sequence of the form
         A*-Aq+M+O.
?2         Modules                                                                  13

This means that M can be generated by 4 elements ol,. . . , wq in such a way
that the module R = {(a,,. .., a,)~A~lCa~o~ = 0) of linear relations
holding between the oi can be generated by p elements.
\end{mathematicalrecord}
\end{proof}


\begin{theorem}[Theorem 2.6]\label{mcrt-auto-theorem-2-6}\dcref{2.6}\source{matsumura-commutative-ring-theory:page-0028}\uses{}
\begin{mathematicalrecord}
Let A be a ring, and suppose that M is an A-module               of finite
presentation. If
           O+K-N-M-0
is an exact sequence and N is finitely generated then so is K.
\end{mathematicalrecord}
\end{theorem}
\begin{proof}
\begin{mathematicalrecord}
By assumption     there exists an exact sequence of the form
L, 2  L, L M -+ 0, where L, and L2 are free modules of finite rank.
From this we get the following commutative diagram (see Appendix B):




   If we write N = At, + ... + A<,,, then there exist ui~L, such that
cp(li)=f(ui).  Set 5: = ti -a(~,); then (p(&)=O, so that we can write
I$ = @(vi) with ~],EK. Let us now prove that
       K=j3(L2)+4,+..++Aq,.
For any ~EK, set $(r]) = xaiti;        then
        $(S - Cailli) = Cai(ti - 5;) = cc(Cuiui)3
and since 0 = (~a(1 aiui) = f(C a,~,.), we can write cuiui = g(u) with UEL,.
Now
        IcIPf?) = crS(u)= 4x Vi) = $411- 1 uiUi),
SOthat q = /I(u) + 1 uiqi, and this proves our assertion.           w
\end{mathematicalrecord}
\end{proof}


\section{Chain conditions}

\begin{theorem}[Theorem 3.1]\label{mcrt-auto-theorem-3-1}\dcref{3.1}\source{matsumura-commutative-ring-theory:page-0030}\uses{}
\begin{mathematicalrecord}
Let A be a ring and M an A-module.
  (i) Let M? c M be a submodule and 9: M + M/M? the natural map. If
N, and N, are submodules of M such that N, c N,, N, n M? = N, n M?
and cp(N,) = cp(N,) then N, = Nz.
   (ii) Let 0 -+ M? -M      + M? + 0 be an exact sequence of A-modules;
if M? and M? are both Noetherian (or both Artinian), then so is M.
   (iii) Let M be a finite A-module; then if A is Noetherian (or Artinian),
so is M.
\end{mathematicalrecord}
\end{theorem}
\begin{proof}
\begin{mathematicalrecord}
(i) is easy, and we leave it to the reader.
   (ii) is obtained by applying (i) to an ascending (respectively descending)
chain of submodules of M.
   (iii) If M is generated by n elements then it is a quotient of the free module
A?, so that it is enough to show that A? is Noetherian (respectively Artinian).
However, this is clear from (ii) by induction on n. w
   For a module M, it is equivalent to say that M has both the a.c.c. and
the d.c.c., or that M has finite length. Indeed, if I(M) < 00 then I(M i) < 1(M,)
for any two distinct submodules MI c M2 c M, so that the two chain
conditions are clear. Conversely, if M has the d.c.c. then we let M, be a
minimal non-zero submodule of M, let M, be a minimal element among all
submodules of M strictly containing M, , and proceed in the same way to
obtain an ascending chain 0 = MO c M 1 c M, c ...; if M also has the a.c.c.
then this chain must stop by arriving at M, so that M has a composition
series,
   Every submodule of the Z-module Z is of the form nZ, so that Z is
Noetherian, but not Artinian. Let p be a prime, and write W for the Z-
module of rational numbers whose denominator is a power of p; then
the Z-module W/Z is not Noetherian, but it is Artinian, since every proper
submodule of W/Z is either 0 or is generated by p-? for n = 1,2,. . . . This
shows that the a.c.c. and d.c.c. for modules are independent conditions,
but this is not the case for rings, as shown by the following result.
16         Commutative        rings and modules
\end{mathematicalrecord}
\end{proof}


\begin{theorem}[Theorem 3.2]\label{mcrt-auto-theorem-3-2}\dcref{3.2}\source{matsumura-commutative-ring-theory:page-0031}\uses{}
\begin{mathematicalrecord}
Y. Akizuki). An Artinian ring is Noetherian.
\end{mathematicalrecord}
\end{theorem}
\begin{proof}
\begin{mathematicalrecord}
Let A be an Artinian ring. It is sufficient to prove that A has finite
length as an A-module. First of all, A has only finitely many maximal ideals.
Indeed,ifp,,p,,...     is an infinite set of distinct maximal ideals then it is easy
 to see that p1 3 plpz =) p1p2p3?. is an infinite descending chain of ideals,
 which contradicts the assumption. Thus, we let pl, pz,. . . , pr be all the
 maximal ideals of A and set I = p1p2.. . p, = rad (A). The descending chain
 III2    2 ?.. stops after finitely many steps, so that there is an s such that
 I? = Z?+i. If we set (0:I?) = J then
           (J:Z) = ((0:P):I) = (O:r+l) = J;
 let?s prove that J = A. By contradiction, suppose that J # A; then there exists
 an ideal J? which is minimal among all ideals strictly bigger than J. For any
 XEJ? - J we have J? = Ax + J. Now I = rad (A) and J # J?, so that by NAK
J? # Ix + J, and hence by minimality of J? we have Ix + J = J, and this gives
 Ix c J. Thus XE(J:I) = J, which is a contradiction. Therefore J = A, so that
 I? = 0. Now consider the chain of ideals



Let M and Mpi be any two consecutive terms in this chain; then MIMpi is a
vector space over the field A/p,, and since it is Artinian, it must be tinite-
dimensional. Hence, l(M/Mp,) < co, and therefore the sum 1(A) of these
terms is also finite. a
\end{mathematicalrecord}
\end{proof}


\begin{theorem}[Theorem 3.3]\label{mcrt-auto-theorem-3-3}\dcref{3.3}\source{matsumura-commutative-ring-theory:page-0031}\uses{}
\begin{mathematicalrecord}
If A is Noetherian then so are A[X] and A[Xj.
\end{mathematicalrecord}
\end{theorem}
\begin{proof}
\begin{mathematicalrecord}
The statement for A[X] is the well-known Hilbert basis theorem
(see, for example Lang, Algebra, or [AM], p. 81), and we omit the proof.
We now briefly run through the proof for A[Xj.              Set B = A[XJ, and
let I be an ideal of B; we will prove that I is finitely generated. Write
Z(r) for the ideal of A formed by the leading coefficients a, of f = a,X? +
a *+1 xr+1 + . . ? as f runs through In X?B; then we have
          I(0) c 1(l) C 1(2) C . . . .
Since A is Noetherian, there is an s such that I(s) = I(s + 1) = ?.; moreover,
each Z(i) is finitely generated. For each i with 0 < i < s we take finitely many
elements aivEA generating Z(i), and choose giVEI nXiB having Ui, as the
coefficient of Xi. These giV now generate 1. Indeed, for f~1 we can take
a linear combination        go of the g,,? with coefficients in A such that
?3       Chain conditions                                                      17

f -geEZn       XB, then take a linear combination   gi of the glV with
coefficients in A such that f - go - g,EZnX?B,    and proceeding in the
same way we get
           f-go-g1     --...-g,ElnX?+?B.
Now Z(s + 1) = Z(s), so we can take a linear combination       gs+ i of the Xg,,
with coefficients in A such that
        f-g90-g1--~~-gs+lEznX?+2B.
We now proceed in the same way to get gs+ *,. . . For i < s, each gi
is a linear combination        of the giy with coefficients in A, and, for
i > s, a combination      of the elements XiPsgsV. For each i > s we write
gi = ~vaivXi-?g,,,     and then for each v we set h,, = ~im,saivXi~S; h, is an
element of B, and
         f=gO+...+gs-l+ChYgSY.             n .

   A ring A[b,,. . .,b,] which is finitely generated as a ring over a
Noetherian ring A is a quotient of a polynomial ring A[X,, . . . ,X,1, and
so by the Hilbert basis theorem is again Noetherian. We now give some
other criteria for a ring to be Noetherian.
\end{mathematicalrecord}
\end{proof}


\begin{theorem}[Theorem 3.4]\label{mcrt-auto-theorem-3-4}\dcref{3.4}\source{matsumura-commutative-ring-theory:page-0032}\uses{}
\begin{mathematicalrecord}
I. S. Cohen). If all the prime ideals of a ring A are finitely
generated then A is Noetherian.
\end{mathematicalrecord}
\end{theorem}
\begin{proof}
\begin{mathematicalrecord}
Write r for the set of ideals of A which are not finitely generated.
If r # fa then by Zorn?s lemma r contains a maximal element I. Then I is
not a prime ideal, so that there are elements x, YEA with x$Z, ~$1 but
xy~l. Now Z+Ay is bigger than I, and hence is finitely generated, so
that we can choose ui,. . .,u,EZ such that
           Z+Ay=(u,      ,..., u,,y).
Moreover,      Z:y= {aEAlayeZ}        contains x, and is thus bigger than
I, so it has a finite system of generators             (ur,. .,u}. Finally, it is
easy to check that Z = (ui ,..., y,u,y,..      .,v,y); hence, Z$r, which is a
contradiction.    Therefore r = 121. n
\end{mathematicalrecord}
\end{proof}


\begin{theorem}[Theorem 3.5]\label{mcrt-auto-theorem-3-5}\dcref{3.5}\source{matsumura-commutative-ring-theory:page-0032}\uses{}
\begin{mathematicalrecord}
Let A be a ring and M an A-module. Then if M is a
Noetherian module, A/ann(M) is a Noetherian ring.
\end{mathematicalrecord}
\end{theorem}
\begin{proof}
\begin{mathematicalrecord}
If we set A = Alann (M) and view M as an A-module, then the
submodules of M as an A-module or A-module coincide, so that M is
also Noetherian as an A-module. We can thus replace A by 2, and then
arm(M) = (0). Now letting M = Ao, + ... + Aw,, we can embed A in M?
by means of the map a++(au r, . . . , aw,). By Theorem 1, M? is a Noetherian
module, so that its submodule A is also Noetherian. (This theorem can
be expressed by saying that a ring having a faithful Noetherian module
is Noetherian.)  n
18       Commutative    rings and modules
\end{mathematicalrecord}
\end{proof}


\begin{theorem}[Theorem 3.6]\label{mcrt-auto-theorem-3-6}\dcref{3.6}\source{matsumura-commutative-ring-theory:page-0033}\uses{}
\begin{mathematicalrecord}
E. Formanek Cl]). Let A be a ring, and B an A-module
which is finitely generated and faithful over A. Assume that the set of
submodules of B of the form 1B with I an ideal of A satisfies the a.c.c.;
then A is Noetherian.
\end{mathematicalrecord}
\end{theorem}
\begin{proof}
\begin{mathematicalrecord}
It will be enough to show that B is a Noetherian A-module. By
contradiction, suppose that it is not; then the set
             IB I is an ideal of A and B/IB is                               /
           i    1non-Noetherian   as A-module   i
contains (0) and so is non-empty, so that by assumption it contains a
maximal element. Let ZB be one such maximal element; then replacing B
by BJIB and A by A/ann(B/IB)         we see that we can assume that B is a
non-Noetherian      A-module, but for any non-zero ideal I of A the quotient
BjZB is Noetherian.
   Next we set
          I = (NIN is a submodule of B and B/N is a faithful A-module}.
If B = Ab, + *a. + Ab, then for a submodule N of B,
         NElYoVaaA-0,          {ab, ,..., ab} $N.
From this, one sees at once that Zorn?s lemma applies to I; hence there
exists a maximal element N, of I. If B/N,, is Noetherian then A is a
Noetherian     ring, and thus B is Noetherian,         which contradicts our
hypothesis. It follows that on replacing B by B/N, we arrive at a module
B with the following properties:
   (1) B is non-Noetherian    as an A-module;
   (2) for any ideal I # (0) of A, B/IB is Noetherian;
   (3) for any submodule N # (0) of B, B/N is not faithful as an A-module.
   Now let N be any non-zero submodule of B. By (3) there is an element
aeA with a # 0 such that a(B/N) = 0, that is such that aB c N. By (2)
B/aB is a Noetherian module, so that N/aB is finitely generated; but since
B is finitely generated so is aB, and hence N itself is finitely generated.
Thus, B is a Noetherian module, which contradicts (1). n
   As a corollary of this theorem we get the following result.
\end{mathematicalrecord}
\end{proof}


\begin{theorem}[Theorem 3.7]\label{mcrt-auto-theorem-3-7}\dcref{3.7}\source{matsumura-commutative-ring-theory:page-0033}\uses{}
\begin{mathematicalrecord}
(i) (Eakin-Nagata    theorem). Let B be a Noetherian ring, and Aa subring
of B such that B is finite over A; then A is also a Noetherian ring.
    (ii) Let B be a non-commutative    ring whose right ideals have the a.c.c.,
and let A be a commutative subring of B. If B is finitely generated as a
left A-module then A is a Noetherian ring.
    (iii) Let B be a non-commutative     ring whose two-sided ideals have the
a.c.c., and let A be a subring contained in the centre of B; if B is finitely
generated as an A-module then A is a Noetherian ring.
93        Chain conditions                                                          19
\end{mathematicalrecord}
\end{theorem}
\begin{proof}
\begin{mathematicalrecord}
B has a unit, so is faithful as an A-module.              Hence it is enough to
apply the previous theorem.      n
\end{mathematicalrecord}
\end{proof}
